\documentclass[11pt,a4paper]{article}
\usepackage[utf8]{inputenc}
\usepackage[T1]{fontenc}
\usepackage{lmodern}
\usepackage[margin=2.3cm]{geometry}
\usepackage{amsmath,amssymb}
\usepackage{graphicx}
\usepackage{booktabs}
\usepackage{float}
\usepackage{xcolor}
\usepackage{listings}
\usepackage{caption}
\usepackage{algorithm}
\usepackage{algpseudocode}
\usepackage[colorlinks=true,linkcolor=blue!50!black,urlcolor=blue!50!black]{hyperref}

% ---- edit these ----
\newcommand{\studentname}{Mesut}
\newcommand{\studentid}{}
\newcommand{\coursename}{DEP 521E -- Special Topics in Earthquake Engineering, Homework 2}

\lstdefinestyle{py}{
  language=Python,
  basicstyle=\ttfamily\footnotesize,
  keywordstyle=\color{blue!70!black}\bfseries,
  commentstyle=\color{green!45!black}\itshape,
  stringstyle=\color{red!60!black},
  numbers=left, numberstyle=\tiny\color{gray}, numbersep=6pt,
  frame=single, rulecolor=\color{gray!60},
  breaklines=true, showstringspaces=false, columns=fullflexible,
  backgroundcolor=\color{gray!4}
}

\title{\textbf{Part C: Bilinear Element Model and Testing\\Coded with AI}}
\author{\studentname\ifx\studentid\empty\else\ (\studentid)\fi\\[2pt]\small \coursename}
\date{\today}

\begin{document}
\maketitle

\begin{abstract}
This report presents a Python implementation of a rate-independent bilinear hysteretic element,
developed with the help of an AI coding assistant as required in Part~C of the assignment.
The element is described by a primary stiffness $k_p$, a secondary stiffness $k_s$ and two
parallel border lines that bound the force. The code is driven through the displacement history
given in the assignment, and the resulting displacement and force time histories and the
force--displacement hysteresis are reported. The results are checked against hand calculations
and against an independent fine-step path-tracing solution.
\end{abstract}

%==============================================================================
\section{Problem definition}
The bilinear element of Fig.~1 in the assignment is defined by the yield displacement $d_y$,
the yield force $f_y$ and the secondary stiffness $k_s$. The primary (elastic) stiffness and the
characteristic strength (the force intercept of the borders) are
\begin{equation}
  k_p = \frac{f_y}{d_y}, \qquad q_y = f_y - k_s d_y .
\end{equation}
At any displacement $x$ the force is bounded by the \emph{top} and \emph{bottom} borders,
\begin{equation}
  F_{\text{top}}(x) = q_y + k_s x, \qquad F_{\text{bot}}(x) = -q_y + k_s x ,
  \label{eq:borders}
\end{equation}
and the region between them is the \emph{interior region}, in which the element responds with
the primary stiffness $k_p$. On a border, loading away from the interior follows the border
with stiffness $k_s$, while unloading moves back into the interior with stiffness $k_p$.

The function required by the assignment has the form
\begin{center}
\texttt{df, param\_a = bilinear(dy, fy, k2, x\_ini, f\_ini, dx, param\_a)}
\end{center}
where \texttt{x\_ini} and \texttt{f\_ini} are the displacement and force at the start of the step,
\texttt{dx} is the displacement increment and \texttt{param\_a} carries the state of the element
(\texttt{interior}, \texttt{top} or \texttt{bottom}). The function returns the force increment
\texttt{df} and the updated state.

%==============================================================================
\section{Formulation}
\subsection{Elastic trial and return to the border}
For a step from $x_i$ to $x_{i+1} = x_i + \Delta x$, an elastic trial force is computed with the
primary stiffness,
\begin{equation}
  F^{\text{trial}} = F_i + k_p\,\Delta x .
\end{equation}
Because $k_s < k_p$, the trial force can lie outside the borders only if the element has yielded
during the step. In that case the step must end on the border that was crossed. Hence
\begin{equation}
  F_{i+1} =
  \begin{cases}
    F_{\text{top}}(x_{i+1}), & F^{\text{trial}} \ge F_{\text{top}}(x_{i+1}) \quad\text{(state: top)},\\[2pt]
    F_{\text{bot}}(x_{i+1}), & F^{\text{trial}} \le F_{\text{bot}}(x_{i+1}) \quad\text{(state: bottom)},\\[2pt]
    F^{\text{trial}},        & \text{otherwise} \quad\text{(state: interior)},
  \end{cases}
  \label{eq:update}
\end{equation}
and $\Delta F = F_{i+1} - F_i$.

This compact rule reproduces every case listed in the hint of the assignment:
\begin{itemize}
  \item \textbf{Interior, no yielding:} the trial stays between the borders and $\Delta F = k_p\Delta x$.
  \item \textbf{Interior, yielding during the step:} the trial passes a border; the force is placed
        on the border, which is the same as moving with $k_p$ up to the border and with $k_s$ afterwards.
  \item \textbf{On the top border, $\Delta x > 0$:} $F^{\text{trial}} = F_{\text{top}}(x_i) + k_p\Delta x
        > F_{\text{top}}(x_{i+1})$, so $\Delta F = k_s\Delta x$ and the state stays \texttt{top}.
  \item \textbf{On the top border, $\Delta x < 0$:} the element unloads with $k_p$ into the interior;
        if the unloading is large enough to reach the bottom border, the step ends on the bottom border.
  \item The bottom border is handled symmetrically.
\end{itemize}
A zero increment returns $\Delta F = 0$ and leaves the state unchanged, which avoids a state
decision based on round-off when the point lies exactly on a border.

The element requires $k_s < k_p$ (otherwise the borders would be steeper than the elastic branch
and the rule above would not hold); the function raises an error if this is violated.

\subsection{Algorithm}
\begin{algorithm}[H]
\caption{Bilinear element, one displacement step}
\begin{algorithmic}[1]
\Require $d_y, f_y, k_s, x_i, F_i, \Delta x$, state
\State $k_p \gets f_y/d_y$, \quad $q_y \gets f_y - k_s d_y$
\If{$\Delta x = 0$} \Return $0$, state \EndIf
\State $x_{i+1} \gets x_i + \Delta x$
\State $F_{\text{top}} \gets q_y + k_s x_{i+1}$, \quad $F_{\text{bot}} \gets -q_y + k_s x_{i+1}$
\State $F^{\text{trial}} \gets F_i + k_p\Delta x$
\If{$F^{\text{trial}} \ge F_{\text{top}}$} \Return $F_{\text{top}} - F_i$, \texttt{top}
\ElsIf{$F^{\text{trial}} \le F_{\text{bot}}$} \Return $F_{\text{bot}} - F_i$, \texttt{bottom}
\Else\ \Return $F^{\text{trial}} - F_i$, \texttt{interior}
\EndIf
\end{algorithmic}
\end{algorithm}

\subsection{Program structure}
The script \texttt{bilinear\_element.py} (Appendix~\ref{app:code}) contains
\begin{itemize}
  \item \texttt{bilinear(...)}: the element function described above;
  \item \texttt{run\_history(x, dy, fy, k2)}: drives the element through a displacement history,
        starting from $F = 0$ in the interior region, and accumulates $F_{i+1} = F_i + \Delta F$;
  \item \texttt{save\_history\_plot}, \texttt{save\_hysteresis\_plot}: the figures;
  \item \texttt{main()}: the test case of the assignment and the printed results table.
\end{itemize}

%==============================================================================
\section{Use of AI}
As required in Part~C, the code was developed with an AI coding assistant (Claude, Anthropic,
used through Claude Code). The assignment text was given to the assistant, which wrote the
element function, the driver, the plotting functions and a separate verification script, ran
them and checked the figures. The formulation in Section~2 and the results in Section~4
were reviewed by the student. The verification described in Section~4.3 was used to confirm
that the AI-generated code is correct rather than relying on the code by inspection.

%==============================================================================
\section{Test of the bilinear element}
\subsection{Input}
The test parameters are listed in Table~\ref{tab:param}. The input displacement history is
\begin{align*}
x = [\,&0,\ 0.5,\ 0.75,\ 1.25,\ 1.5,\ 1.0,\ 1.75,\ 0.75,\ 0.25,\ 0.5,\ 0.0,\ -0.5,\\
       &-0.75,\ -1.5,\ -1.0,\ -1.75,\ -2.0,\ -2.5,\ -2.25,\ -1.75,\ -1.0,\ -0.25,\ 0,\ 1.0\,].
\end{align*}
No time step is given, so the time axis of the time histories is the step number.

\begin{table}[H]
\centering
\caption{Input and derived parameters.}
\label{tab:param}
\begin{tabular}{lcc}
\toprule
Parameter & Symbol & Value\\
\midrule
Yield displacement      & $d_y$ & 1\\
Yield force             & $f_y$ & 5\\
Secondary stiffness     & $k_s$ & 0.25\\
\midrule
Primary stiffness       & $k_p = f_y/d_y$       & 5\\
Characteristic strength & $q_y = f_y - k_s d_y$ & 4.75\\
\bottomrule
\end{tabular}
\end{table}

\subsection{Results}
The computed forces and states are listed in Table~\ref{tab:results}. The displacement and force
time histories are shown in Figs.~\ref{fig:disp} and~\ref{fig:force}, and the hysteresis loop
with the two borders in Fig.~\ref{fig:hyst}.

\begin{table}[H]
\centering
\caption{Force and state at each step.}
\label{tab:results}
\small
\begin{tabular}{rrrl@{\hspace{2.2em}}rrrl}
\toprule
Step & $x$ & $F$ & State & Step & $x$ & $F$ & State\\
\midrule
 0 &  0.00 &  0.0000 & interior & 12 & $-0.75$ & $-4.9375$ & bottom\\
 1 &  0.50 &  2.5000 & interior & 13 & $-1.50$ & $-5.1250$ & bottom\\
 2 &  0.75 &  3.7500 & interior & 14 & $-1.00$ & $-2.6250$ & interior\\
 3 &  1.25 &  5.0625 & top      & 15 & $-1.75$ & $-5.1875$ & bottom\\
 4 &  1.50 &  5.1250 & top      & 16 & $-2.00$ & $-5.2500$ & bottom\\
 5 &  1.00 &  2.6250 & interior & 17 & $-2.50$ & $-5.3750$ & bottom\\
 6 &  1.75 &  5.1875 & top      & 18 & $-2.25$ & $-4.1250$ & interior\\
 7 &  0.75 &  0.1875 & interior & 19 & $-1.75$ & $-1.6250$ & interior\\
 8 &  0.25 & $-2.3125$ & interior & 20 & $-1.00$ &  2.1250 & interior\\
 9 &  0.50 & $-1.0625$ & interior & 21 & $-0.25$ &  4.6875 & top\\
10 &  0.00 & $-3.5625$ & interior & 22 &  0.00 &  4.7500 & top\\
11 & $-0.50$ & $-4.8750$ & bottom & 23 &  1.00 &  5.0000 & top\\
\bottomrule
\end{tabular}
\end{table}

\begin{figure}[H]
  \centering
  \includegraphics[width=0.86\textwidth]{displacement_history.png}
  \caption{Displacement time history.}
  \label{fig:disp}
\end{figure}

\begin{figure}[H]
  \centering
  \includegraphics[width=0.86\textwidth]{force_history.png}
  \caption{Force time history.}
  \label{fig:force}
\end{figure}

\begin{figure}[H]
  \centering
  \includegraphics[width=0.78\textwidth]{hysteresis.png}
  \caption{Hysteresis: force vs.\ displacement. Dashed lines are the top and bottom borders,
           Eq.~\eqref{eq:borders}.}
  \label{fig:hyst}
\end{figure}

\paragraph{Discussion.}
From step~0 to step~2 the element is elastic with slope $k_p = 5$. Yielding occurs during step~3
(at $x = 1$, $F = 5$) and the response then follows the top border with slope $k_s = 0.25$.
Unloading in step~5 and step~7 follows the elastic slope; in step~7 the displacement decreases by
1.0, which changes the force by $-5$ and leaves the point inside the interior region because the
elastic range between the borders is $2q_y = 9.5$ in force. Reloading in step~6 reaches the top
border again. In the negative direction the element yields on the bottom border during step~11
and reaches its minimum force $F = -5.375$ at the peak displacement $x = -2.5$. Reloading from
step~18 to step~21 crosses the whole interior region with slope $k_p$ and reaches the top border
during step~21. The hysteresis loop of Fig.~\ref{fig:hyst} is a parallelogram whose elastic
branches are parallel to each other and whose plastic branches lie on the two borders, which is
the expected bilinear behaviour.

\subsection{Verification}
The AI-generated code was checked in three independent ways:
\begin{enumerate}
  \item \textbf{Hand calculation.} Steps 1--11 were computed by hand. For example, in step~3 the
        trial force is $3.75 + 5(0.5) = 6.25 > F_{\text{top}}(1.25) = 4.75 + 0.25(1.25) = 5.0625$,
        so $F = 5.0625$; in step~11 the trial force is $-3.5625 - 5(0.5) = -6.0625 <
        F_{\text{bot}}(-0.5) = -4.875$, so $F = -4.875$. All hand values agree exactly with
        Table~\ref{tab:results}.
  \item \textbf{Independent path tracing.} Each step was divided into 20\,000 sub-steps and the
        force was integrated with the incremental stiffness ($k_s$ when loading outward on a border,
        $k_p$ otherwise). The largest difference from the bilinear function is about
        $8\times10^{-12}$.
  \item \textbf{Consistency checks.} At every step the force lies between the two borders, and
        every step reported as \texttt{top} or \texttt{bottom} lies exactly on that border.
\end{enumerate}

%==============================================================================
\section{Conclusions}
\begin{itemize}
  \item A bilinear element was implemented in Python with the signature required by the
        assignment, using an elastic trial and a return to the crossed border.
  \item The element was tested with the given parameters ($d_y = 1$, $f_y = 5$, $k_s = 0.25$)
        and displacement history. The maximum and minimum forces are $5.1875$ (at $x = 1.75$)
        and $-5.375$ (at $x = -2.5$).
  \item The results agree with hand calculations and with an independent fine-step solution,
        so the AI-generated code is considered correct for this test.
\end{itemize}

%==============================================================================
\appendix
\section{Python code}\label{app:code}
\lstinputlisting[style=py]{bilinear_element.py}

\end{document}
